\begin{essay}{}
	
ПОДДЕРЖКА ПРИНЯТИЯ РЕШЕНИЙ, ПОШАГОВАЯ ИНСТРУКЦИЯ, ГРАФОВОЕ ПРЕДСТАВЛЕНИЕ, МЕТОД АНАЛИЗА ИЕРАРХИЙ (МАИ), ДЕТЕРМИНИРОВАННЫЙ АЛГОРИТМ, РЕКОМЕНДАЦИИ, АВАРИЙНЫЕ СЦЕНАРИИ, PYTHON

Объектом исследования являются методы поддержки принятия решений при выполнении пошаговых инструкций, заданных в графовом формате.

Объектом разработки является программное обеспечение, реализующее метод поддержки принятия решений на основе графового представления инструкции и метода анализа иерархий с поддержкой аварийных сценариев и режима экспертной отладки.

Целью работы является разработка метода поддержки принятия решений при выполнении пошаговой инструкции, заданной в графовом формате.

В аналитическом разделе выполнен анализ предметной области систем поддержки принятия решений, сопровождающих выполнение набора предписаний человеком, проведён сравнительный анализ методов принятия решений, описан существующий графовый формат описания инструкций, на примере выбранного домена формализованы классы возможных состояний, событий и связанных с ними действий, выполнена формализация постановки задачи в виде функциональной модели.

В конструкторском разделе разработан метод поддержки принятия решений в рамках выполнения пошаговой инструкции, заданной в графовом формате, описаны основные особенности предлагаемого метода, сформулированы ограничения предметной области, изложены ключевые этапы метода в виде функциональной модели и схем алгоритмов, выделены структуры данных, необходимые для реализации метода.

В технологическом разделе обоснован выбор средств программной реализации, разработано программное обеспечение, реализующее представленный метод, выполнено его тестирование, описано взаимодействие пользователя с программным обеспечением.

В исследовательском разделе описана исследовательская процедура, составлен набор сценариев развития событий для каждой решаемой задачи в рамках выбранного домена, включая различные виды нештатного развития ситуации, выполнена параметризация метода согласно описанной исследовательской процедуре, даны рекомендации о применимости метода.

Результаты проведённого исследования подтверждают начальные предположения о том, что предложенная модификация МАИ с группировкой альтернатив позволяет существенно сократить временные затраты эксперта по сравнению с классическим подходом, а экспертная настройка весов критериев в режиме отладки повышает качество выдаваемых рекомендаций.
\end{essay}